% !TEX root = ../main.tex
\chapter{Lender/Dealer of Last Resort}\label{ch:lolr}
\begin{paracol}{2}

\begin{sources}
\begin{tabularx}{\linewidth}{@{}l l L@{}}
\toprule
\textbf{Segment} & \textbf{Length} & \textbf{Content}\\
\midrule
\seg{12}{1}{L12.1 FT: Citibank and the SIVs} & 5:10 & Pandit resigns; Citi's structured investment vehicles and the collapse of ABCP.\\
\seg{12}{2}{L12.2 The Art of Central Banking} & 3:42 & Hawtrey (1932); the Fed's crisis facilities.\\
\seg{12}{3}{L12.3 Evolution of Monetary Policy: 1951--1987} & 7:06 & Quantity and price views; operating targets from free reserves to Fed funds.\\
\seg{12}{4}{L12.4 The Taylor Rule: 1987--2007} & 7:04 & The rule, the Fisher effect, inflation targeting.\\
\seg{12}{5}{L12.5 Monetary Transmission Mechanism} & 10:00 & From Fed funds to term rates to asset prices, through three dealer markets.\\
\seg{12}{6}{L12.6 Anatomy of a Normal Crisis} & 8:18 & Dealers and banks absorb an imbalance, for a price.\\
\seg{12}{7}{L12.7 Anatomy of a Serious Crisis} & 4:32 & The Fed as backstop of the banks that backstop the dealers.\\
\seg{12}{8}{L12.8 Should the Fed Intervene or Not?} & 8:40 & Elasticity against spirals, discipline against accumulated mistakes.\\
\seg{12}{9}{L12.9 The Fed as Dealer of Last Resort: 2007--2009} & 15:50 & The stages of the crisis in balance sheets.\\
\bottomrule
\end{tabularx}

\smallskip
\textbf{Files.} Transcripts \file{materials/transcripts/L12-P*.txt}.
\textbf{Readings.} Stigum, chapter~9 on monetary policy; R.~G.~Hawtrey, \emph{The Art of Central Banking} (1932), chapter~4 (quoted).
\end{sources}

\section*{Motivation}
The last lecture before the midterm puts the pieces together: the three dealer markets of lectures 6--11 (Fed funds, term funding, bonds) form the
\emph{transmission mechanism} of monetary policy; normal fluctuations are absorbed by dealers and banks; serious crises need the Fed as lender of last
resort to the banks; and in 2007--2009 the Fed went further, becoming the money market itself and finally a \emph{dealer of last resort}. The lecture is
mostly about lender of last resort in the money market; dealer of last resort in capital markets needs derivatives and comes after the midterm
\ts{12}{2}{0:07}.

% ------------------------------------------------------------------
\section{FT: Citibank and the SIVs}\label{sec:lr-ft}

\begin{casestudy}[FT, October 2012: Vikram Pandit resigns]
Citigroup's CEO Vikram Pandit resigned under board pressure and was replaced by Michael Corbat, who had run Citi Holdings, where the toxic assets were put to
be worked out \ts{12}{1}{0:07}\ts{12}{1}{0:28}\ts{12}{1}{3:31}. Citi, built by Sandy Weill as a ``state-of-the-art financial supermarket'', came close to
nationalisation in the crisis \ts{12}{1}{0:49}\ts{12}{1}{1:09}. The FT's timeline omits an episode Mehrling would add: in November 2007, as Citi announced losses
that grew every few days (5, then 8, then 11 billion), it received a \$7.5 billion capital infusion from the Abu Dhabi Investment Authority: professional
investors still thought it was a minor blip, as many economists thought subprime too small to matter \ts{12}{1}{1:30}\ts{12}{1}{1:50}\ts{12}{1}{2:10}\ts{12}{1}{2:51}.
It was the beginning of the global financial crisis \ts{12}{1}{3:11}.
\end{casestudy}

\paragraph{The SIVs.} Citi was a famous part of the shadow banking system: it set up off-balance-sheet \emph{structured investment vehicles} that held AAA
tranches of residential mortgage-backed securities, funded with asset-backed commercial paper \ts{12}{1}{3:31}\ts{12}{1}{3:52}\ts{12}{1}{4:12}. The first thing that
happened in the crisis was the collapse of the SIVs and of the ABCP market: 90-day paper funding 30-year mortgages; at maturity the holders, money funds,
declined to roll, the SIV found no other funding and borrowed from its host, Citibank, which ended up with the assets on its own balance sheet
\ts{12}{1}{4:33}\ts{12}{1}{4:53}.

\begin{definition}[New concepts: structured investment vehicle, ABCP]
A \emph{structured investment vehicle} (SIV) is an off-balance-sheet entity sponsored by a bank that holds long-term assets (here AAA mortgage-backed securities)
funded by short-term debt. \emph{Asset-backed commercial paper} (ABCP) is commercial paper secured by a pool of assets, typically 30 to 90 days.
\end{definition}

% ------------------------------------------------------------------
\section{The art of central banking}\label{sec:lr-art}

Ralph Hawtrey, \emph{The Art of Central Banking} (London, 1932, ``an auspicious year''), chapter 4, subtitled ``the lender of last resort'': ``a central bank is
a bankers' bank'', and the central bank is the lender of last resort; Hawtrey, in the line of Bagehot and up through Keynes, Sayers and Goodhart, was an
inspiration for Mehrling \ts{12}{2}{0:48}\ts{12}{2}{1:08}\ts{12}{2}{1:29}. (His copy, withdrawn from the Kirkland House library at Harvard, cost three dollars
\ts{12}{2}{1:50}.)

\paragraph{The crisis facilities.} The New York Fed lists the Fed's lending facilities: regular open market operations and the discount window, and then a long
list that was not there before the crisis \ts{12}{2}{2:10}\ts{12}{2}{2:31}: the term discount window program, the Term Auction Facility, the Primary Dealer Credit
Facility, transitional credit extensions, reciprocal currency arrangements (the liquidity swaps), securities lending and the Term Securities Lending Facility, the
ABCP Money Market Mutual Fund Liquidity Facility, the Commercial Paper Funding Facility, the Money Market Investor Funding Facility, the Term Asset-Backed
Securities Loan Facility \ts{12}{2}{2:51}\ts{12}{2}{3:12}. ``This is what put a floor under it'' \ts{12}{2}{3:33}.

% ------------------------------------------------------------------
\section{Normal monetary policy}\label{sec:lr-normal}

\subsection{Quantity or price?}
Once central banks accepted Bagehot's responsibility, they started to think how to avoid crises in the first place, and invented monetary policy, which evolved from
1873 to what it was before the crisis \ts{12}{3}{0:07}\ts{12}{3}{0:28}. Stigum shows an ambivalence about what the Fed controls: a quantity (money, reserves, a balance
sheet measure) or a price (the Fed funds rate)? \ts{12}{3}{0:48}\ts{12}{3}{1:09} Those who emphasise quantities have the equation of exchange $MV=PT$ in mind; those who
emphasise rates, the money market view \ts{12}{3}{1:30}\ts{12}{3}{1:51}. The two can be translated into each other:
\begin{itemize}
  \item \textbf{quantity}: the Fed sits on top of the hierarchy (the pyramid of lecture~2) and controls the best money of all, its own liability, high-powered money,
        which it uses to influence the balance between discipline and elasticity \ts{12}{3}{2:12}\ts{12}{3}{2:33}\ts{12}{3}{2:54}\ts{12}{3}{3:15};
  \item \textbf{price}: the Fed pegs the shortest end of the term structure, an administered policy rate, while the market determines the long end, security prices
        \ts{12}{3}{3:38}\ts{12}{3}{3:59}.
\end{itemize}
Either way, monetary policy is about the \emph{transmission mechanism} from what the Fed controls, its own balance sheet, to the financial system and the economy
\ts{12}{3}{4:20}.

\paragraph{Operating targets} (Stigum) \ts{12}{3}{4:41}\ts{12}{3}{5:02}\ts{12}{3}{5:24}:
\begin{itemize}
  \item 1951: free reserves;
  \item 1979: non-borrowed reserves (Volcker's ``Saturday night special'');
  \item 1983: borrowed reserves;
  \item 1987: the Fed funds rate, ever since.
\end{itemize}
Mehrling is not sure the Fed changed much what it actually did, rather how it explained it \ts{12}{3}{5:45}. In 1979, facing double-digit inflation, Volcker presented
himself as a convert to monetarism; most people think it was a convenient language: saying ``I'm going to jack the Fed funds rate up to 15\%'' might have cost him his
confirmation, saying ``I'm a monetarist'' had the same consequence without saying it, and he stopped inflation in its tracks \ts{12}{3}{6:05}\ts{12}{3}{6:26}\ts{12}{3}{6:46}.

\begin{history}[The Saturday night special]
On Saturday 6 October 1979 the FOMC, chaired by Paul Volcker since August, announced that it would control the supply of bank reserves (non-borrowed reserves) instead of
the Fed funds rate; the rate rose to around 20\% in 1980--1981. The approach was relaxed in 1982.\par
\emph{References:} Federal Reserve press release, 6 Oct 1979; A.~H.~Meltzer, \emph{A History of the Federal Reserve}, vol.~2, Chicago UP, 2009.
\end{history}

\subsection{The Taylor rule (1987--2007)}
1987 is when Mehrling got his job; by the eve of the crisis almost all economists agreed on what the Fed was doing and should do: the \emph{Taylor rule}
\ts{12}{4}{0:07}\ts{12}{4}{0:28}.

\begin{keyformulas}
\textbf{Taylor rule} (in the form given in class) \ts{12}{4}{0:49}:
\[
  r = \rho + \pi^e + \alpha\,(\pi^e - \pi^\ast) + \beta\,(y - y_f),
\]
with $r$ the target (Fed funds) rate, $\rho$ the natural real rate of interest (the return on real capital), $\pi^e$ expected inflation, $\pi^\ast$ the inflation target,
$y-y_f$ the gap between output and its full-employment level \ts{12}{4}{1:09}\ts{12}{4}{1:30}\ts{12}{4}{1:51}.

\medskip
\textbf{Fisher effect}: $r = \rho + \pi^e$, nominal rate $=$ real rate $+$ expected inflation \ts{12}{4}{2:31}\ts{12}{4}{2:56}.
\end{keyformulas}

\begin{itemize}
  \item As a regression, the rule fits the behaviour of central banks around the world well in normal times (not in the crisis): a positive description; a large literature
        then made it normative \ts{12}{4}{1:51}\ts{12}{4}{2:11}\ts{12}{4}{2:31}.
  \item The first part, $\rho+\pi^e$, is the \emph{Fisher effect} (Irving Fisher, Yale): what markets do by themselves, compensating creditors for expected inflation. Fisher
        himself thought people often got inflation wrong, causing fluctuations; modern economics, with rational expectations, assumes $\pi^e$ is unbiased
        \ts{12}{4}{2:56}\ts{12}{4}{3:17}\ts{12}{4}{3:37}\ts{12}{4}{3:58}.
  \item The rest is the Fed leaning against the wind: if expected inflation exceeds the target, raise the rate by more than the difference (more discipline); if output is
        below full employment, lower it (more elasticity) \ts{12}{4}{4:18}\ts{12}{4}{4:39}.
\end{itemize}
Careers were made on the optimal $\alpha$ and $\beta$, on whether to add a financial-stability term: ``inflation targeting'' spread all over the world. Economics came to
believe it had mastered the economy to the point of arguing whether a coefficient should be 1.1 or 1.2; the Great Moderation, then the crisis. ``Nobody believes this
anymore'' \ts{12}{4}{4:59}\ts{12}{4}{5:19}\ts{12}{4}{5:40}\ts{12}{4}{6:01}. Mehrling, who was reading Hawtrey instead, was better prepared: ``in the land of the blind the
one-eyed man is king'' \ts{12}{4}{6:22}\ts{12}{4}{6:42}.

\begin{coursenote}[Taylor's original rule]
John Taylor (``Discretion versus policy rules in practice'', \emph{Carnegie--Rochester Conference Series}, 1993) wrote $i = \pi + 0.5\,y + 0.5\,(\pi-2) + 2$, with
$\pi$ inflation over the previous four quarters, $y$ the percentage output gap and 2\% for both the real rate and the inflation target. The total response of $i$ to
inflation is $1+0.5=1.5>1$ (the ``Taylor principle''). In the form above the coefficient on $\pi^e$ is $1+\alpha$, so the principle requires $\alpha>0$; in class
Mehrling says ``$\alpha$ is greater than one'' \ts{12}{4}{4:18}.
\end{coursenote}

% ------------------------------------------------------------------
\section{The monetary transmission mechanism}\label{sec:lr-transmission}

The three dealer diagrams of lectures 10--11, in reverse order, form the transmission mechanism \ts{12}{5}{0:08}\ts{12}{5}{0:28}:
\begin{enumerate}
  \item \textbf{Fed funds} (overnight): the Fed makes the outside spread (IOER, discount rate); the horizontal axis is \emph{settlement risk}, the risk of pushing payments to
        tomorrow \ts{12}{5}{1:11}\ts{12}{5}{1:32}\ts{12}{5}{1:53}\ts{12}{5}{2:16};
  \item \textbf{term funding} (three months): \emph{liquidity risk}, borrowing short and lending long; upward-sloping quotes in yields \ts{12}{5}{2:37}\ts{12}{5}{2:58};
  \item \textbf{bonds}: asset prices, downward-sloping quotes; \emph{price risk} \ts{12}{5}{2:58}\ts{12}{5}{3:19}.
\end{enumerate}

\begin{nfigure}
\begin{tikzpicture}[font=\scriptsize,x=0.64cm,y=0.64cm]
  \foreach \xs/\t/\xl/\yl in {0/{Fed funds (overnight)}/{settlement risk}/{rate},5.6/{term funding}/{liquidity risk}/{term rate},11.2/{bonds}/{price risk}/{price}} {
    \begin{scope}[xshift=\xs cm]
      \draw[-{Stealth[length=3pt]}] (0,0) -- (5.4,0) node[below left]{\xl};
      \draw[-{Stealth[length=3pt]}] (0,0) -- (0,4.4) node[above]{\yl};
      \node[font=\scriptsize\bfseries] at (2.7,5.2) {\t};
    \end{scope}}
  \begin{scope}[xshift=0cm]
    \draw[thick,thmcol] (0,3.6) -- (5.2,3.6) node[right]{discount};
    \draw[thick,thmcol] (0,0.6) -- (5.2,0.6) node[right]{IOER};
    \draw[thick,defcol] (0.5,1.2) -- (4.8,3.0);
    \draw[dashed,keycol] (0,2.1) -- (5.2,2.1) node[right]{target};
  \end{scope}
  \begin{scope}[xshift=5.6cm]
    \draw[thick,defcol] (0.4,1.6) -- (4.9,3.6);
    \draw[dashed,keycol] (0,2.1) -- (5.2,2.1);
  \end{scope}
  \begin{scope}[xshift=11.2cm]
    \draw[thick,defcol] (0.4,3.6) -- (4.9,1.4);
  \end{scope}
  \draw[-{Stealth[length=5pt]},very thick,keycol] (6.9,2.8) -- (7.8,2.8);
  \draw[-{Stealth[length=5pt]},very thick,keycol] (14.9,2.8) -- (15.8,2.8);
\end{tikzpicture}
\caption{The transmission mechanism as three dealer markets \ts{12}{5}{0:28}\ts{12}{5}{3:43}: the Fed sets the overnight rate within its outside spread; the term rate must exceed it
by enough to pay dealers for liquidity risk; bond prices depend on the cost of funding dealers' inventories. Raising the Fed funds target raises term rates and lowers bond prices.}
\label{fig:lr-transmission}
\end{nfigure}

The story \ts{12}{5}{4:04}\ts{12}{5}{6:29}\ts{12}{5}{6:49}\ts{12}{5}{7:11}:
\begin{enumerate}
  \item the Fed sets a target Fed funds rate and enforces it by trading (temporary open market operations); the level of the target is a choice between leaning toward
        discipline and leaning toward elasticity, and the same rate can mean elasticity in some circumstances and discipline in others \ts{12}{5}{4:25}\ts{12}{5}{4:46}\ts{12}{5}{5:07};
  \item raising the target (and with it the discount rate) narrows the gap between Fed funds and the term rate, which dealers do not like: they raise the term rate;
  \item funding becomes more expensive, so bond dealers want smaller inventories: downward pressure on bond prices;
  \item asset prices, mortgage-backed securities among them, affect the price of loans and so the real economy.
\end{enumerate}
This works through the price of money and the price of capital, and instantly, through arbitrage in securities markets; it does not need banks lending multiples of their
reserves (the money multiplier) or a bank-lending channel: institutions that no longer exist \ts{12}{5}{7:31}\ts{12}{5}{7:51}\ts{12}{5}{8:12}\ts{12}{5}{8:33}. The Fed may not move prices
exactly as it wants, because the market watches and anticipates it; but ``monetary policy does matter'' in a deregulated system. What it \emph{should} do has to be rethought
\ts{12}{5}{8:53}\ts{12}{5}{9:14}\ts{12}{5}{9:34}.

% ------------------------------------------------------------------
\section{Anatomy of a normal crisis}\label{sec:lr-normalcrisis}

A ``normal crisis'' is the kind of imbalance that can happen every day \ts{12}{6}{0:07}. Households (non-banks) decide, for whatever reason, maybe only for a day, that they
want to hold money instead of securities \ts{12}{6}{0:28}\ts{12}{6}{0:49}. The dealer system exists to absorb imbalances in order flow \ts{12}{6}{1:09}:
\begin{taccts}
\tacct[2.0cm]{Households}{$-$Securities\\$+$Deposits}{}\quad
\tacct[2.0cm]{Dealers}{$+$Securities}{$+$Repo (from bank)}\quad
\tacct[2.0cm]{Banks}{$+$Repo (to dealer)}{$+$Deposits}
\end{taccts}
The household sector sells securities on net; dealers buy them and finance them with repo from banks, which lend by expanding their balance sheets
\ts{12}{6}{1:32}\ts{12}{6}{1:53}\ts{12}{6}{2:14}\ts{12}{6}{2:34}. ``Sort of magical'': households think they sold to other households, but the securities were absorbed by the dealer and
banking system, and their deposits are \emph{new} deposits, indistinguishable from old ones \ts{12}{6}{2:55}\ts{12}{6}{3:16}.

\paragraph{The price of absorption.} Dealers and banks do it to make money \ts{12}{6}{3:36}. The dealer, seeing the imbalance in its order book, buys at a lower price; when
flows reverse, it sells higher. The bank lends at a slightly higher rate, and later liquidates at a lower one \ts{12}{6}{3:57}\ts{12}{6}{4:17}\ts{12}{6}{4:38}\ts{12}{6}{4:58}\ts{12}{6}{5:20}.
The households pay without seeing it, through a lower price for their bonds \ts{12}{6}{5:40}\ts{12}{6}{6:00}\ts{12}{6}{6:21}. These price distortions away from fundamental value are
possibly \emph{equilibrating}: low prices attract other buyers, the positions are liquidated \ts{12}{6}{6:43}\ts{12}{6}{7:03}. Even a seller with no information, who just needs
cash for a payment, moves the price, because the buyer fears he knows something; it comes back once the market realises he does not \ts{12}{6}{7:23}\ts{12}{6}{7:43}\ts{12}{6}{8:04}.

\section{Anatomy of a serious crisis}\label{sec:lr-seriouscrisis}

In a normal crisis there is no Fed: the private system absorbs it and the Fed funds rate stays put \ts{12}{7}{0:07}. A big shock makes dealers and banks anxious; that is
when the Fed comes in \ts{12}{7}{0:28}\ts{12}{7}{0:48}. As in Bagehot's world, the central bank sets a rate above the market rate; if the market on its own would push rates above
it, the central bank becomes lender of last resort \ts{12}{7}{1:10}\ts{12}{7}{1:32}:
\begin{taccts}
\tacct[2.0cm]{Fed}{$+$Loans to banks}{$+$Reserves}\quad
\tacct[2.0cm]{Banks}{$+$Repo to dealers\\$+$Reserves}{$+$Deposits\\$+$Borrowing from Fed}\quad
\tacct[2.0cm]{Dealers}{$+$Securities}{$+$Repo}
\end{taccts}
If the Fed backstops the banks, the banks are more willing to backstop the dealers; this keeps a large shock within bounds \ts{12}{7}{1:52}\ts{12}{7}{2:23}\ts{12}{7}{2:43}.

\begin{proposition}[Why the central bank can stop a crisis]
In a crisis everyone wants money and nobody wants securities; the securities must go somewhere. The lender of last resort works by letting the unwanted securities be absorbed
into the financial sector and creating the cash people want to fund them: expansion of dealers, of banks, ultimately of the Fed. The Fed backstops the banks' short positions
in money, i.e.\ their liquidity risk, and it cannot itself run into trouble with short positions in money, because its liability is the ultimate money, at least domestically:
it operates as a money market dealer \ts{12}{7}{3:04}\ts{12}{7}{3:24}\ts{12}{7}{3:45}\ts{12}{7}{4:06}\ts{12}{7}{4:26}.
\end{proposition}

% ------------------------------------------------------------------
\section{Should the Fed intervene or not?}\label{sec:lr-should}

Private dealers absorb fluctuations for profit; the Fed's motive is different: when should it take the pressure off, and when let it ride? The balance between discipline and
elasticity \ts{12}{8}{0:09}\ts{12}{8}{0:29}\ts{12}{8}{0:50}. The Fed can always provide elasticity, but that does not mean it should \ts{12}{8}{1:11}.

\paragraph{For elasticity.} Prices may become so distorted that people make bad decisions; or distortions may stop being equilibrating: falling asset prices make assets
unusable as collateral, holders dump them, a \emph{liquidity spiral}, and the system collapses \ts{12}{8}{1:11}\ts{12}{8}{1:32}\ts{12}{8}{4:16}\ts{12}{8}{4:36}.

\paragraph{The value investors.} A student asks: why would value-based investors not step in? In this picture the Fed backstops banks, not dealers (dealer of last resort in
the money market) \ts{12}{8}{2:13}\ts{12}{8}{2:33}. It may intervene directly in asset markets, as it did, because the value investors did not: they refused, or lost their shirts,
or wait for prices of 20 cents on the dollar, which would come only with a great depression; the Fed is not there to make money for value investors, it is concerned with
unemployment. If value investors do step in, the Fed need not \ts{12}{8}{2:54}\ts{12}{8}{3:14}\ts{12}{8}{3:34}\ts{12}{8}{3:55}.

\paragraph{For discipline.} Why not always be elastic? Minsky: what moves money market rates is the mismatch between cash commitments and cash flows; if rates never moved,
nothing would stop agents from pushing commitments into the future forever. Discipline makes mistakes get recognised: at some point you pay or go bankrupt
\ts{12}{8}{4:56}\ts{12}{8}{5:16}\ts{12}{8}{5:37}\ts{12}{8}{5:57}. Otherwise mistakes build up like dry tinder, and ``Smokey the Bear'' produces larger forest fires
\ts{12}{8}{6:18}. In stress people are forced to sell at prices they would never accept, and so to recognise their mistakes and start afresh; but too much discipline ``kills the
patient'' \ts{12}{8}{6:59}\ts{12}{8}{7:20}.

\begin{intuition}[The art of central banking]
Central bankers ``feel this in their bones'': not every hiccup, not even every significant stress, deserves elasticity. In this crisis the Fed did not flood the market at once:
it let some bad mortgages ``walk'' while preventing collapse; then things got worse than expected and it did more, and more \ts{12}{8}{6:38}\ts{12}{8}{7:40}\ts{12}{8}{8:00}\ts{12}{8}{8:21}.
\end{intuition}

% ------------------------------------------------------------------
\section{The Fed as dealer of last resort, 2007--2009}\label{sec:lr-dolr}

The Fed's balance sheet in the crisis (lecture~3) shows the stages \ts{12}{9}{0:07}\ts{12}{9}{0:28}:
\begin{enumerate}
  \item \textbf{Interest rates} (August 2007 to Bear Stearns): Fed funds from 5\% to 2\%, to make it more profitable for banks to take on liquidity risk and support asset prices
        indirectly \ts{12}{9}{0:49}\ts{12}{9}{1:09};
  \item \textbf{composition} (after Bear Stearns): of about a trillion of Treasury bills, \$500 billion sold and the proceeds lent to brokers and dealers \ts{12}{9}{1:09}\ts{12}{9}{1:30};
  \item \textbf{expansion} (after Lehman and AIG): about two trillion more of deposits and of lending; the interbank market broke down and ``the Fed became the interbank market'':
        surplus institutions deposited at the Fed, deficit institutions borrowed from it \ts{12}{9}{1:30}\ts{12}{9}{1:51};
  \item \textbf{capital market}: purchases of mortgage-backed securities \ts{12}{9}{2:11}.
\end{enumerate}

\paragraph{The crisis in balance sheets} \ts{12}{9}{2:32}--\ts{12}{9}{10:15}:
\begin{enumerate}
  \item \textbf{Before}: Citi's SIV holds RMBS funded by ABCP, held by a money fund that issues shares to (international) investors; nothing is on Citi's or the Fed's balance sheet.
  \item \textbf{First stage}, which looked like a normal crisis: the money fund declines to roll the ABCP; Citi replaces the funding with a loan to the SIV, funded on its own
        balance sheet, say with repo. The fund now holds a liability of Citibank with recourse to its whole balance sheet. Citi did this as a business: take assets off customers who
        no longer want them, and charge for it \ts{12}{9}{3:56}\ts{12}{9}{4:17}\ts{12}{9}{4:38}\ts{12}{9}{4:58}\ts{12}{9}{5:19}.
  \item \textbf{Second stage}: the fund prefers Citi's unsecured commercial paper, based on its whole book (Lehman too issued such paper, held by the Reserve Primary Fund), or
        Eurodollar deposits at foreign banks \ts{12}{9}{5:39}\ts{12}{9}{6:00}\ts{12}{9}{6:20}\ts{12}{9}{7:00}. The shadow banking system collapses onto the traditional banking system,
        backstopped by the Fed and the FDIC \ts{12}{9}{7:23}\ts{12}{9}{7:43}.
  \item \textbf{Third stage}: European banks, with little access to the Fed, cannot roll their dollar funding when money funds stop renewing Eurodollar deposits
        \ts{12}{9}{8:04}\ts{12}{9}{8:25}. The Fed lends dollars to foreign central banks through liquidity swaps (about \$600 billion at the peak), and the Treasury issues bills
        through a supplementary account at the Fed, so that money funds can hold what they want, Treasury bills \ts{12}{9}{8:48}\ts{12}{9}{9:10}\ts{12}{9}{9:34}\ts{12}{9}{9:55}\ts{12}{9}{10:15}.
\end{enumerate}
\begin{taccts}
\tacct[2.0cm]{Fed $+$ Treasury}{$+$Swap (loan to ECB)}{$+$Treasury bills}\quad
\tacct[2.0cm]{ECB}{$+$Dollar loans to European banks}{$+$Swap (dollars owed to Fed)}\quad
\tacct[2.0cm]{Money fund}{$-$Eurodollar deposits\\$+$Treasury bills}{}
\end{taccts}

\paragraph{From the best money down the hierarchy.} The sequence: first the normal absorption, letting prices move to give the banking system an incentive; then lower Fed funds;
then direct entry into the term funding market (liquidity swaps, lending to dealers, the term facilities), backstopping banks; finally, as value investor, the mortgage-backed
securities market \ts{12}{9}{10:35}\ts{12}{9}{10:56}\ts{12}{9}{11:16}\ts{12}{9}{11:37}. The Fed moved down the hierarchy, putting more and more of the financial system on its own
balance sheet. At no point did it want to: central bankers are conservative, they do not want credit risk or interest rate risk, which is why their money is the best money; but in
a crisis, ``lend freely at a high rate against what would be good security in a normal time'' \ts{12}{9}{11:57}\ts{12}{9}{12:18}.

\begin{intuition}[A floor under the system]
Normally the Fed acts at a distance, setting the overnight rate and letting profit-seeking dealers and arbitrage transmit it \ts{12}{9}{12:39}\ts{12}{9}{12:59}. When imbalances built up
over years are too large, the crisis is not temporary and keeps getting worse, and the Fed's responsibility as ultimate backstop forces it to do more, until it has put a floor under
the system. It can, because its liabilities are the best money: ultimately it could buy every risky security for cash, though that would be the end of the market system
\ts{12}{9}{13:19}\ts{12}{9}{13:40}\ts{12}{9}{14:01}\ts{12}{9}{14:21}. It always hopes to do just enough; the ECB is doing that now, hoping the pressure will make politicians act: perhaps
``a slow-motion collapse'' of the European system, in which long periods of nothing alternate with bursts of excitement \ts{12}{9}{14:41}\ts{12}{9}{15:02}\ts{12}{9}{15:23}\ts{12}{9}{15:43}.
\end{intuition}

\begin{definition}[New concepts: dealer of last resort]
A central bank acts as \emph{dealer of last resort} when, instead of only lending to banks against collateral (lender of last resort), it stands ready to buy and sell, or to
backstop the prices of, securities that private dealers will no longer absorb, putting a floor under their prices.
\end{definition}

% ------------------------------------------------------------------
\flushside
\section*{Key formulas and ideas}
\begin{keyformulas}
\begin{tabularx}{\linewidth}{@{}l L@{}}
Quantity vs price & controlling the best money (top of the hierarchy) $\equiv$ pegging the short end of the term structure\\
Taylor rule & $r=\rho+\pi^e+\alpha(\pi^e-\pi^\ast)+\beta(y-y_f)$; Fisher effect $r=\rho+\pi^e$\\
Transmission & Fed funds (settlement risk) $\to$ term rate (liquidity risk) $\to$ bond prices (price risk) $\to$ economy\\
Normal crisis & dealers and banks absorb order imbalances for a price (temporary distortions)\\
Serious crisis & Fed backstops banks, banks backstop dealers; the Fed cannot fail on short positions in its own money\\
Intervene? & elasticity against liquidity spirals; discipline to make mistakes recognised\\
2007--2009 & rates $\to$ composition $\to$ expansion (Fed $=$ interbank market) $\to$ swaps, facilities $\to$ MBS: dealer of last resort\\
\end{tabularx}
\end{keyformulas}

\section*{Exercises}

\begin{exercise}[Taylor rule]
With $\rho=2\%$, $\pi^\ast=2\%$, $\alpha=0.5$, $\beta=0.5$:
\begin{enumerate}[(a)]
  \item compute the prescribed rate for $\pi^e=3\%$ and an output gap of $-2\%$;
  \item what rate does the rule prescribe for $\pi^e=1\%$ and a gap of $-6\%$? What happens at the zero lower bound?
\end{enumerate}
\end{exercise}

\begin{exercise}[Transmission]
Using the three diagrams of \cref{fig:lr-transmission}, explain step by step what happens to the three-month rate and to bond prices when the Fed raises its target by 25 basis points,
and why the effect may be smaller if the market had already anticipated it.
\end{exercise}

\begin{exercise}[A normal crisis]
Households sell \$100 of bonds. Dealers buy them at a price 1\% below fundamental value, financed by repo at the term rate. Three days later the flow reverses and dealers sell at
fundamental value. Write the T-accounts at each step and compute the dealers' profit if repo costs 0.5\% a year.
\end{exercise}

\begin{exercise}[The SIV unwinds]
An SIV holds RMBS of 100 funded by ABCP of 100 held by a money fund. Write the T-accounts (SIV, Citi, money fund, Fed) for each stage of \cref{sec:lr-dolr}, ending with the money fund
holding Treasury bills.
\end{exercise}
